\documentclass[11pt]{article}

\usepackage[T1]{fontenc}
\usepackage{lmodern}
\usepackage[margin=1.05in]{geometry}
\usepackage{microtype}
\usepackage{amsmath,amssymb,amsthm,mathtools}
\usepackage{enumitem}
\usepackage[hidelinks]{hyperref}

\setlist{leftmargin=*,itemsep=0.25em,topsep=0.45em}
\setlength{\parindent}{1.2em}
\setlength{\parskip}{0.12em}
\setlength{\emergencystretch}{2em}
\allowdisplaybreaks
\numberwithin{equation}{section}

\newtheorem{theorem}{Theorem}[section]
\newtheorem{proposition}[theorem]{Proposition}
\newtheorem{lemma}[theorem]{Lemma}
\newtheorem{corollary}[theorem]{Corollary}
\theoremstyle{definition}
\newtheorem{definition}[theorem]{Definition}
\theoremstyle{remark}
\newtheorem{remark}[theorem]{Remark}

\newcommand{\C}{\mathbb C}
\newcommand{\N}{\mathbb N}
\newcommand{\Pp}{\mathcal P}
\newcommand{\M}{\mathcal M}
\newcommand{\A}{\mathcal A}
\newcommand{\D}{\mathcal D}
\newcommand{\dd}{\mathrm d}
\newcommand{\ord}{\operatorname{ord}}
\newcommand{\Frac}{\operatorname{Frac}}
\newcommand{\trdeg}{\operatorname{trdeg}}

\title{\bfseries A Differential-Algebraic Obstruction to Polynomial-Gauge\\
Painlev\'e III Representations of the Riemann Xi Function}
\author{Michael Haufschild}
\date{}

\begin{document}
\maketitle

\begin{abstract}
Let
\[
 \xi(s)=\frac12s(s-1)\pi^{-s/2}\Gamma(s/2)\zeta(s)
\]
be the completed Riemann xi function.  A recent bridge conjecture proposes that a fixed-monodromy
Painlev\'e~III function of type $D_6$, pulled back by $t=s(1-s)$ and modified by the exponential
of a polynomial of degree at most two, can equal $\xi$.  We prove that no such representation
exists.  In fact, the polynomial degree, the two monodromy exponents, and the local monodromy data
are unrestricted.

The obstruction is differential algebraic.  Translation by two of the logarithmic derivative of
$\xi$ exposes the Dirichlet series with coefficients $\Lambda(n)(1-n^{-2})$.  Finite
prime-coordinate specializations and a Vandermonde determinant make every finite family of its
differential jets algebraically independent; a polynomial-weighted uniqueness theorem then
transfers this independence from Dirichlet coefficients to analytic functions.  This gives a
direct proof that $\xi$ is hypertranscendental over $\C(s)$.  On the Painlev\'e side, elimination
of the sigma function without division by the tau function produces an explicit nonzero
third-order algebraic differential equation for every regular local $\mathrm{PIII}(D_6)$ tau
function.  Differential algebraicity survives the quadratic pullback and every polynomial
exponential gauge.  Equality with $\xi$ is therefore impossible.  No hypothesis about the zeros
of $\xi$, including the Riemann hypothesis, enters the argument.
\end{abstract}

\noindent\textbf{2020 Mathematics Subject Classification.}
Primary 12H05, 34M55; Secondary 11M06, 11J91, 34M40.

\noindent\textbf{Keywords.}
Riemann xi function; Painlev\'e III; isomonodromic tau function; differential algebra;
hypertranscendence; Dirichlet series.

\section{Introduction}
\label{sec:introduction}

Algebraic differential equations provide a structural invariant for proposed exact
representations of special functions.  A function is differentially algebraic over $\C(s)$ when
it and finitely many derivatives satisfy a polynomial relation with coefficients in $\C(s)$.
This class is closed under field operations, differentiation, algebraic changes of variable, and
exponentials of polynomials.  Consequently, any function constructed from a fixed solution of a
finite-order algebraic ordinary differential equation by these operations remains
differentially algebraic.  Hypertranscendental functions, which satisfy no such relation, cannot
be represented in that way.  The point of this invariant is that it tests an exact bridge before
questions of zero location, asymptotic accuracy, or parameter fitting arise.

For $\theta_0,\theta_\infty\in\C$, the sigma form of Painlev\'e~III of type $D_6$ considered
here is
\begin{equation}
 (t\sigma'')^2
 =\bigl(2\sigma'-\theta_0-\theta_\infty\bigr)^2
      (\sigma-t\sigma')
   +4\sigma'(\sigma'-\theta_0)(\sigma'-\theta_\infty).
 \label{eq:sigma-form-intro}
\end{equation}
The associated tau function $u$ is related to $\sigma$ by
\begin{equation}
 \sigma(t)=t\frac{u'(t)}{u(t)}.
 \label{eq:tau-sigma-quotient}
\end{equation}
This is the standard Hamiltonian tau-function description of the $D_6$ equation
\cite{JMU1981,Miwa1981,Okamoto1981,Okamoto1987}.  In a 2026 preprint, Planat wrote the same sigma
equation as his equation~(5), proposed the quadratic pullback
$t=s(1-s)$ and the carrier
\begin{equation}
 E_{D_6}(s)=e^{-2s}u\bigl(s(1-s)\bigr).
 \label{eq:planat-carrier}
\end{equation}
This is his equation~(18).  His Conjecture~3 asks whether the monodromy data can be selected so
that
\begin{equation}
 \xi(s)=e^{Q(s)}E_{D_6}(s),
 \qquad Q\in\C[s],\quad \deg Q\leq2.
 \label{eq:planat-bridge}
\end{equation}
The constant-gauge form was excluded in the same work, while the polynomial-gauge form and the
selection of monodromy data were left open \cite[Conjecture~3 and Section~7]{Planat2026}.

Equation \eqref{eq:sigma-form-intro} already determines the differential-algebraic class of every
solution satisfying this fixed-parameter sigma equation.  What requires proof is that the
conclusion remains usable after the quadratic pullback, at possible zeros of the tau function,
and after passage from a local tau chart to an alleged global entire bridge.  The core elimination
below uses the cross-multiplied identity
\begin{equation}
 u\sigma=tu',
 \label{eq:tau-sigma-cross}
\end{equation}
so no division by $u$ enters the proof.

\subsection*{Relation to previous work}

Differential transcendence of the Riemann zeta function belongs to a classical line beginning
with Ostrowski's theory of Dirichlet series and algebraic differential equations
\cite{Ostrowski1920}; Reich developed broad hypertranscendence criteria for Dirichlet series
\cite{Reich1984}, and Nagoshi gives a modern account and extensions to automorphic
$L$-functions \cite{Nagoshi2016}.  Joint differential-algebraic relations between $\zeta$ and
$\Gamma$ have also been studied extensively, from Markus's partial independence results
\cite{Markus2007} through later theorems of Han--Liu and Wang--Yao
\cite{HanLiu2021,WangYao2024}.  Those results impose strong differential-independence constraints
on combinations of $\zeta$ and $\Gamma$.  We therefore do not claim that
hypertranscendence of completed xi is new in isolation.

The contribution here is a proof adapted to the proposed bridge.  It isolates an exact
translation defect of $\xi'/\xi$ whose coefficients have a nonzero pivot at every prime, proves
the required differential independence directly in the arithmetic-function ring, and combines
that result with an explicit denominator-free elimination of the $D_6$ tau equation.  The outcome
is a degree-unrestricted obstruction for the particular Painlev\'e architecture.  The argument
uses neither value-distribution theory nor a joint-independence theorem for $\zeta$ and $\Gamma$.

The other side of the obstruction is the following unconditional statement.

\begin{theorem}[Hypertranscendence of completed xi]
\label{thm:xi-ht-intro}
The function $\xi$ satisfies no nonzero algebraic differential equation over $\C(s)$.
\end{theorem}

We prove this in Sections~\ref{sec:dirichlet} and~\ref{sec:xi}.  Its arithmetic
input is rigid: the $+2$ shift of $\xi'/\xi$, after subtraction of an explicit
rational function, is
\begin{equation}
 \sum_{n\geq1}\Lambda(n)(1-n^{-2})n^{-s}.
 \label{eq:prime-source-intro}
\end{equation}
At every prime $p$, the coefficient is $\log p\,(1-p^{-2})\neq0$.  Distinct prime logarithms
produce a Vandermonde minor for every finite differential jet.  Polynomial-weighted uniqueness
of Dirichlet series then shows that \eqref{eq:prime-source-intro} is hypertranscendental over the
actual coordinate field $\C(s)$.

Our main result is stronger than the conjectured degree-two exclusion.  A \emph{regular local
$D_6$ tau chart} means a nonempty open set on which $u$ and $\sigma$ are holomorphic and satisfy
\eqref{eq:tau-sigma-cross} and \eqref{eq:sigma-form-intro}.  Every ordinary nonzero
isomonodromic tau function has such a chart around each point outside its zero and singular
sets; local analyticity of isomonodromic tau functions is established in
\cite{JMU1981,Miwa1981}.

\begin{theorem}[Polynomial-gauge $\mathrm{PIII}(D_6)$ no-go]
\label{thm:main-intro}
Let $\theta_0,\theta_\infty\in\C$ and $Q\in\C[s]$.  There do not exist a function
$E:\C\to\C$, a disc $U\subset\C$ whose closure avoids $s=\tfrac12$, and a regular local $D_6$ tau chart
$(u,\sigma)$ on $g(U)$, where $g(s)=s(1-s)$, such that
\begin{align}
 E(s)&=e^{-2s}u(g(s)) &&(s\in U),\label{eq:local-carrier}\\
 \xi(s)&=e^{Q(s)}E(s) &&(s\in\C).\label{eq:global-bridge}
\end{align}
Consequently, \eqref{eq:planat-bridge} is false.  The conclusion holds without any restriction
on $\deg Q$.
\end{theorem}

The proof uses neither the location nor the density of the zeros of $\xi$.  It also does not
compare large-$t$ asymptotics of Painlev\'e solutions with the Riemann--von Mangoldt formula.
It proves the narrower structural fact needed here: every tau function satisfying the fixed
$\mathrm{PIII}(D_6)$ sigma equation lies locally in the differential-algebraic class, whereas
$\xi$ does not.

Section~\ref{sec:da} establishes the differential-algebraic closure and pullback results used
later.  Section~\ref{sec:dirichlet} proves the coefficient and analytic independence theorems.
Section~\ref{sec:xi} proves Theorem~\ref{thm:xi-ht-intro} from the exact shift identity.
Section~\ref{sec:tau} derives the denominator-free tau equation, and
Section~\ref{sec:nogo} proves Theorem~\ref{thm:main-intro} and its degree-two corollary.

\section{Differential-algebraic preliminaries}
\label{sec:da}

Let $K$ be a differential field of characteristic zero, with derivation denoted by $\partial$.
For an element $f$ of a differential field extension, write
\[
 K\langle f\rangle=K(f,\partial f,\partial^2f,\ldots).
\]

\begin{definition}
An element $f$ is \emph{differentially algebraic over $K$} if there are $m\geq0$ and a nonzero
polynomial $P\in K[Y_0,\ldots,Y_m]$ such that
\[
 P(f,\partial f,\ldots,\partial^m f)=0.
\]
It is \emph{hypertranscendental over $K$} otherwise.
\end{definition}

For meromorphic functions, we take $K=\C(s)$ and $\partial=\dd/\dd s$.  Multiplying through by a
common denominator shows that the definition is unchanged if one requires polynomial rather than
rational coefficients in $s$.  These are the standard notions of differential algebra
\cite{Ritt1950,Kolchin1973}.

\begin{proposition}[Differential-field closure]
\label{prop:da-field}
In any differential field extension of $K$, the elements differentially algebraic over $K$ form
a differential subfield.  In particular, if $f$ and $h$ are differentially algebraic, then so are
$f+h$, $fh$, $\partial f$, and $f/h$ when $h\neq0$.
\end{proposition}

\begin{proof}
We recall the algebraic argument.  If
$P(f,\ldots,\partial^m f)=0$, choose such a relation with $m$ minimal and then with minimal degree
in $\partial^m f$.  It follows that $\partial^m f$ is algebraic over
$K(f,\ldots,\partial^{m-1}f)$.  In characteristic zero, differentiation of its minimal polynomial
expresses $\partial^{m+1}f$ algebraically over the same field enlarged by $\partial^m f$: the
coefficient of $\partial^{m+1}f$ is the nonzero formal derivative of the separable minimal
polynomial.  Induction gives
\[
 \trdeg_K K\langle f\rangle<\infty.
\]
Conversely, finite transcendence degree of $K\langle f\rangle$ forces a finite set of the jets of
$f$ to be algebraically dependent, so it yields a differential equation.

If $f$ and $h$ are differentially algebraic over $K$, then $h$ remains differentially algebraic
over the larger base $K\langle f\rangle$.  Applying the preceding finite-transcendence-degree
criterion twice shows that $K\langle f,h\rangle$ has finite transcendence degree over $K$.
Every element of this differential field, including the four expressions in the statement, is
therefore differentially algebraic over $K$.
\end{proof}

\begin{lemma}[Polynomial exponential gauges]
\label{lem:exp-gauge}
For every $Q\in\C[s]$, the function $e^{Q(s)}$ is differentially algebraic over $\C(s)$ and is
nowhere zero.  Multiplication by $e^{Q(s)}$ preserves differential algebraicity and
hypertranscendence.
\end{lemma}

\begin{proof}
The function $v=e^Q$ satisfies the nonzero first-order equation
\[
 v'-Q'v=0.
\]
It is nowhere zero because $e^z e^{-z}=1$.  Proposition~\ref{prop:da-field} gives closure under
multiplication.  If $e^Qf$ were differentially algebraic, multiplication by the differentially
algebraic inverse $e^{-Q}$ would make $f$ differentially algebraic.  This proves the converse and
hence preservation of both classes.
\end{proof}

\begin{lemma}[Translation]
\label{lem:translation}
If $f$ is differentially algebraic over $\C(s)$, then $s\mapsto f(s+a)$ is differentially
algebraic over $\C(s)$ for every $a\in\C$.
\end{lemma}

\begin{proof}
Replace $s$ by $s+a$ in a differential equation for $f$.  Translation is an automorphism of
$\C(s)$ and commutes with $\dd/\dd s$, so the resulting differential polynomial is nonzero and
annihilates $f(s+a)$.
\end{proof}

We next record the pullback result in the precise local form needed below.

\begin{proposition}[Nonconstant algebraic pullback]
\label{prop:pullback}
Let $g\in\C[s]$ be nonconstant, let $U$ be a nonempty open set on which $g'$ has no zero, and let
$u$ be holomorphic on $g(U)$.  If $u$ satisfies a nonzero algebraic differential equation over
$\C(t)$ on $g(U)$, then $y=u\circ g$ satisfies a nonzero algebraic differential equation over
$\C(s)$ on $U$.
\end{proposition}

\begin{proof}
Work first in differential-polynomial rings.  The substitution $t\mapsto g(s)$ embeds
$\C(t)$ into $\C(s)$ because $g$ is nonconstant.  Define
\begin{equation}
 Z_0=Y_0,
 \qquad
 Z_{k+1}=\frac{1}{g'(s)}\frac{\dd Z_k}{\dd s},
 \label{eq:pullback-recursion}
\end{equation}
where $\dd Y_j/\dd s=Y_{j+1}$.  Induction gives the triangular formula
\begin{equation}
 Z_k=(g')^{-k}Y_k+R_k(s,Y_0,\ldots,Y_{k-1}),
 \qquad R_k\in\C(s)[Y_0,\ldots,Y_{k-1}].
 \label{eq:triangular-jets}
\end{equation}
For each finite $m$, the change of variables
$(Y_0,\ldots,Y_m)\mapsto(Z_0,\ldots,Z_m)$ is triangular with nonzero diagonal.  Its inverse is
constructed recursively: $Y_0=Z_0$, and once $Y_0,\ldots,Y_{k-1}$ have been expressed in the
$Z$-variables, equation~\eqref{eq:triangular-jets} gives
\[
 Y_k=(g')^k\bigl(Z_k-R_k(s,Y_0,\ldots,Y_{k-1})\bigr).
\]
Thus substitution by the $Z_k$ is an automorphism of
$\C(s)[Y_0,\ldots,Y_m]$.  Extension of scalars along the injective map
$\C(t)\hookrightarrow\C(s)$ is itself injective.  Consequently, the homomorphism
\[
 \Phi:\C(t)\{Y\}\longrightarrow\C(s)\{Y\},
 \qquad t\mapsto g(s),\quad \partial_t^kY\mapsto Z_k,
\]
is injective.

If $P(t;u,u',\ldots,u^{(m)})=0$, the chain rule and
\eqref{eq:pullback-recursion} give
\[
 Z_k\bigl(s;y,y',\ldots,y^{(k)}\bigr)=u^{(k)}(g(s)).
\]
Thus $\Phi(P)$ annihilates $y$ on $U$.  Injectivity gives $\Phi(P)\neq0$; clearing its finitely
many rational denominators produces the required nonzero equation over $\C(s)$.
\end{proof}

\begin{lemma}[Local identity and global differential algebraicity]
\label{lem:local-global}
Let $F$ be entire on $\C$.  If a nonzero differential polynomial with coefficients in $\C(s)$
annihilates $F$ on a nonempty open set, then $F$ is differentially algebraic over $\C(s)$ on all
of $\C$.
\end{lemma}

\begin{proof}
Clear coefficient denominators.  The resulting expression
\[
 H(s)=P\bigl(s;F(s),F'(s),\ldots,F^{(m)}(s)\bigr)
\]
is entire and vanishes on a nonempty open subset after at most finitely many points have been
removed.  The identity theorem gives $H\equiv0$.  The cleared
differential polynomial remains nonzero because multiplication by a nonzero common denominator is
injective in the integral domain $\C(s)\{Y\}$.
\end{proof}

\section{Prime-supported Dirichlet series and differential independence}
\label{sec:dirichlet}

Let $\A$ be the $\C$-algebra of arithmetic functions $a:\N_{\geq1}\to\C$, with Dirichlet
convolution
\[
 (a*b)(n)=\sum_{d\mid n}a(d)b(n/d).
\]
The sum is finite.  Associativity follows by indexing both $(a*b)*c$ and $a*(b*c)$ by triples
$d_1d_2d_3=n$; the identity is $\varepsilon(1)=1$ and $\varepsilon(n)=0$ for $n>1$.

\begin{lemma}[Arithmetic-function domain]
\label{lem:arithmetic-domain}
The convolution algebra $\A$ is an integral domain.
\end{lemma}

\begin{proof}
Let $a,b\neq0$, and let $m$ and $n$ be the least positive integers for which $a(m)\neq0$ and
$b(n)\neq0$, respectively.  In the divisor sum for $(a*b)(mn)$, a nonzero term indexed by $d$
would require $d\geq m$ and $mn/d\geq n$.  The second inequality gives $d\leq m$, so $d=m$ is
the only possible nonzero term.  Hence
\[
 (a*b)(mn)=a(m)b(n)\neq0.
\]
Thus $a*b\neq0$.  This also supplies directly the domain assertion classically obtained from the
formal-series model of arithmetic functions \cite{CashwellEverett1959}.
\end{proof}

Define the logarithmic derivation
\begin{equation}
 (\delta a)(n)=-(\log n)a(n).
 \label{eq:log-derivation}
\end{equation}
Indeed, $\log n=\log d+\log(n/d)$ for $d\mid n$, and hence
\[
 \delta(a*b)=(\delta a)*b+a*(\delta b).
\]
For $\Re s$ sufficiently large, write
\[
 \D(a)(s)=\sum_{n\geq1}a(n)n^{-s}.
\]
Absolute convergence permits rearrangement and termwise differentiation, giving
\begin{equation}
 \D(a*b)=\D(a)\D(b),
 \qquad
 \frac{\dd}{\dd s}\D(a)=\D(\delta a).
 \label{eq:dirichlet-hom}
\end{equation}

The first lemma isolates the prime-coordinate mechanism.

\begin{lemma}[Prime-pivot differential independence]
\label{lem:prime-pivot}
Suppose that $a\in\A$ is supported on positive prime powers, $a(1)=0$, and $a(p)\neq0$ for every
prime $p$.  Then the infinite family
\[
 a,\delta a,\delta^2a,\ldots
\]
is algebraically independent over $\C$ in the fraction field of $\A$.
\end{lemma}

\begin{proof}
We make precise the formal-series algebra used in the argument.  Let
$M=\N_0^{(\Pp)}$ be the free commutative monoid of finitely supported functions
$\alpha:\Pp\to\N_0$.  Define $\C[[M]]$ to be the set of all coefficient functions
$F:M\to\C$, with pointwise addition and convolution
\[
 (FG)(\alpha)=\sum_{\beta+\gamma=\alpha}F(\beta)G(\gamma).
\]
For fixed $\alpha$, the sum has only
$\prod_{p\in\operatorname{supp}\alpha}(\alpha(p)+1)$ terms.  Therefore multiplication is
well-defined even though infinitely many variables are available and no finiteness condition is
imposed on the total support of $F$.  We write its elements suggestively as
$\sum_{\alpha\in M}F(\alpha)X^\alpha$.

Unique factorization gives a monoid isomorphism
$n\mapsto(v_p(n))_{p\in\Pp}$ from positive integers under multiplication to $M$.  It induces the
ring isomorphism
\begin{equation}
 \Theta(a)=\sum_{n\geq1}a(n)\prod_{p\in\Pp}X_p^{v_p(n)}.
 \label{eq:arithmetic-series}
\end{equation}
Indeed, the coefficient of the monomial corresponding to $n$ in a product is precisely the
finite divisor sum defining Dirichlet convolution.  Lemma~\ref{lem:arithmetic-domain} justifies
the fraction field used in the statement.

Fix $m\geq0$ and choose distinct primes $p_0,\ldots,p_m$.  Put
$F_j=\Theta(\delta^j a)$.  Since $a(1)=0$, every $F_j$ has zero constant term.  After setting all
prime variables other than $X_{p_0},\ldots,X_{p_m}$ equal to zero, we obtain an ordinary formal
power series in $m+1$ variables.  This specialization is a ring homomorphism: a product monomial
uses no excluded prime exactly when both of its factors use no excluded prime.  The degree-one
part of the specialized $F_j$ is
\begin{equation}
 L_j=\sum_{i=0}^{m}a(p_i)(-\log p_i)^jX_{p_i}.
 \label{eq:linear-prime-part}
\end{equation}
The coefficient matrix of the linear forms $L_0,\ldots,L_m$ has determinant
\begin{equation}
 \det\bigl(a(p_i)(-\log p_i)^j\bigr)_{0\leq i,j\leq m}
 =\left(\prod_{i=0}^{m}a(p_i)\right)
   \prod_{0\leq i<k\leq m}\bigl(\log p_i-\log p_k\bigr),
 \label{eq:vandermonde}
\end{equation}
up to the harmless sign determined by the ordering.  It is nonzero: every prime coefficient is
nonzero and the real logarithm is strictly increasing on the positive reals.  Thus the $L_j$ are
linearly independent.

Suppose that a nonzero polynomial $P\in\C[Y_0,\ldots,Y_m]$ satisfied
$P(F_0,\ldots,F_m)=0$.  Let $P_d$ be the nonzero homogeneous part of $P$ of least total degree.
Because every $F_j$ has zero constant term, the degree-$d$ homogeneous part of
$P(F_0,\ldots,F_m)$ is exactly $P_d(L_0,\ldots,L_m)$.  The linear substitution
$Y_j\mapsto L_j$ is injective by \eqref{eq:vandermonde}, so this homogeneous part is nonzero, a
contradiction.  Hence every finite initial family is algebraically independent.  The embedding of
$\A$ into its fraction field preserves this independence.
\end{proof}

To pass from coefficient independence to independence over the coordinate $s$, ordinary
uniqueness of Dirichlet series is not quite enough: differential equations have rational
functions of $s$ as coefficients.  The needed strengthened form follows from the same
least-frequency argument.

\begin{lemma}[Polynomial-weighted Dirichlet uniqueness]
\label{lem:weighted-uniqueness}
Let $b_0,\ldots,b_d$ be arithmetic functions whose Dirichlet series converge absolutely in a
common right half-plane.  If
\begin{equation}
 \sum_{r=0}^{d}s^r\D(b_r)(s)=0
 \label{eq:weighted-relation}
\end{equation}
there, then $b_r(n)=0$ for every $r$ and every $n\geq1$.
\end{lemma}

\begin{proof}
Assume otherwise, and let $n_0$ be the least positive integer for which
\[
 B_{n_0}(X)=\sum_{r=0}^{d}b_r(n_0)X^r
\]
is nonzero.  Choose a real $x_0$ in the common half-plane of absolute convergence.  Evaluating
\eqref{eq:weighted-relation} on the real ray and multiplying by $n_0^x$ gives
\begin{equation}
 B_{n_0}(x)=-\sum_{n>n_0}B_n(x)(n_0/n)^x,
 \qquad B_n(X)=\sum_{r=0}^{d}b_r(n)X^r.
 \label{eq:least-frequency}
\end{equation}
Set $\rho=n_0/(n_0+1)<1$.  For $x\geq\max\{x_0,1\}$ and $n>n_0$,
\[
 (n_0/n)^x
 =n_0^{x_0}n^{-x_0}(n_0/n)^{x-x_0}
 \leq n_0^{x_0}n^{-x_0}\rho^{x-x_0}.
\]
Moreover,
$|B_n(x)|\leq(1+x^d)\sum_{r=0}^d|b_r(n)|$ for $x\geq1$.  Consequently, absolute convergence at
$x_0$ gives a constant $C>0$ such that
\begin{equation}
 |B_{n_0}(x)|\leq C(1+x^d)\rho^x.
 \label{eq:exponential-tail}
\end{equation}
The right-hand side tends to zero exponentially.  A nonzero polynomial cannot satisfy this
bound: if its degree is $k$ and its leading coefficient is $c\neq0$, then
$|B_{n_0}(x)|\geq |c|x^k/2$ for all sufficiently large $x$.  This contradicts
\eqref{eq:exponential-tail}.  Hence no such $n_0$ exists, proving the assertion.  This argument is
the polynomial-weighted version of the classical uniqueness theorem for Dirichlet series
\cite{HardyRiesz1915}.
\end{proof}

\begin{theorem}[Analytic transfer]
\label{thm:analytic-transfer}
Let $a$ satisfy the hypotheses of Lemma~\ref{lem:prime-pivot}, and suppose that $\D(a)$ has a
finite abscissa of absolute convergence.  Then $\D(a)$ is hypertranscendental over $\C(s)$ in its
half-plane of convergence.
\end{theorem}

\begin{proof}
Suppose that $F=\D(a)$ satisfies a nonzero relation over $\C(s)$.  Clear denominators and write it
as
\begin{equation}
 R(s;F,F',\ldots,F^{(m)})=0,
 \qquad R\in\C[s,Y_0,\ldots,Y_m]\setminus\{0\}.
 \label{eq:analytic-diff-relation}
\end{equation}
Expand $R$ in the coordinate:
\[
 R(s;Y)=\sum_{r=0}^{d}s^rP_r(Y),
 \qquad P_r\in\C[Y_0,\ldots,Y_m].
\]
By \eqref{eq:dirichlet-hom}, on a sufficiently far right half-plane,
\[
 P_r(F,F',\ldots,F^{(m)})
 =\D\bigl(P_r(a,\delta a,\ldots,\delta^m a)\bigr).
\]
We verify the convergence assertion used here.  If $\D(a)$ converges absolutely at $x_0$ and
$j\geq0$, then for every $\epsilon>0$ there is $C_{j,\epsilon}>0$ with
$(\log n)^j\leq C_{j,\epsilon}n^\epsilon$ for $n\geq1$; this follows by maximizing
$x^je^{-\epsilon x}$ for $x=\log n$.  Hence
\[
 \sum_{n\geq1}|(\delta^ja)(n)|n^{-(x_0+\epsilon)}
 \leq C_{j,\epsilon}\sum_{n\geq1}|a(n)|n^{-x_0}<\infty.
\]
If $b$ and $c$ converge absolutely at $x$, the finite-divisor rearrangement and Tonelli's theorem
give
\[
 \sum_{n\geq1}|(b*c)(n)|n^{-x}
 \leq
 \left(\sum_{d\geq1}|b(d)|d^{-x}\right)
 \left(\sum_{e\geq1}|c(e)|e^{-x}\right).
\]
Thus every finite sum and convolution occurring in the displayed $P_r$ has a finite abscissa of
absolute convergence.  Lemma~\ref{lem:weighted-uniqueness} applied to
\eqref{eq:analytic-diff-relation} forces
\[
 P_r(a,\delta a,\ldots,\delta^m a)=0
 \quad\text{for every }r.
\]
Lemma~\ref{lem:prime-pivot} then gives $P_r=0$ for every $r$, contradicting $R\neq0$.
\end{proof}

\section{Hypertranscendence of the completed xi function}
\label{sec:xi}

We use the von Mangoldt function
\[
 \Lambda(n)=
 \begin{cases}
  \log p,&n=p^k\text{ for a prime }p\text{ and }k\geq1,\\
  0,&\text{otherwise}.
 \end{cases}
\]
Define
\begin{equation}
 c(n)=\Lambda(n)(1-n^{-2}),
 \qquad
 \mathcal F(s)=\sum_{n\geq1}c(n)n^{-s}.
 \label{eq:filtered-vm}
\end{equation}

\begin{proposition}
\label{prop:filtered-ht}
The series $\mathcal F$ converges absolutely for $\Re s>1$ and is hypertranscendental over
$\C(s)$.
\end{proposition}

\begin{proof}
The bound $0\leq\Lambda(n)\leq\log n$ follows immediately from the definition: in the nonzero
case $n=p^k$ and $\log p\leq\log n$.  Also $|1-n^{-2}|\leq1$ for $n\geq1$.  The series
$\sum_{n\geq2}(\log n)n^{-x}$ converges for $x>1$, as follows from the integral test applied to
$(\log t)t^{-x}$.  Thus \eqref{eq:filtered-vm} converges absolutely in the stated half-plane.

Its coefficient function is supported on positive prime powers and $c(1)=0$.  At a prime $p$,
\[
 c(p)=\log p\,(1-p^{-2})\neq0.
\]
The hypotheses of Lemma~\ref{lem:prime-pivot} and
Theorem~\ref{thm:analytic-transfer} are therefore satisfied.
\end{proof}

\begin{lemma}[Continuation and a zero-free half-plane]
\label{lem:xi-analytic}
The function initially defined for $\Re s>1$ by
\begin{equation}
 \xi(s)=\frac12s(s-1)\pi^{-s/2}\Gamma(s/2)\zeta(s)
 \label{eq:xi-def}
\end{equation}
extends to an entire function.  It has no zero in $\Re s>1$.
\end{lemma}

\begin{proof}
We give the analytic continuation because the local-to-global step later uses entireness.  For
$x>0$, put
\[
 \vartheta(x)=\sum_{n\in\mathbb Z}e^{-\pi n^2x},
 \qquad
 \psi(x)=\sum_{n\geq1}e^{-\pi n^2x}.
\]
The Fourier transform, with kernel $e^{-2\pi ixy}$, of $e^{-\pi x y^2}$ is
$x^{-1/2}e^{-\pi \nu^2/x}$.  One direct verification differentiates the transform with respect to
$\nu$, integrates by parts, and solves the resulting first-order equation; its value at $\nu=0$
is the Gaussian integral.  Periodizing the Gaussian gives an absolutely and locally uniformly
convergent Fourier series.  Its $k$th Fourier coefficient is the transform at $k$, so evaluation
at the origin yields
\begin{equation}
 \vartheta(x)=x^{-1/2}\vartheta(1/x).
 \label{eq:theta-transform}
\end{equation}
This is Poisson summation for the Gaussian, derived here in the only form needed below.

For $\Re s>1$, absolute convergence permits termwise integration, and the change of variable
$y=\pi n^2x$ gives
\begin{align*}
 \int_0^\infty\psi(x)x^{s/2}\frac{\dd x}{x}
 &=\sum_{n\geq1}\int_0^\infty e^{-\pi n^2x}x^{s/2}\frac{\dd x}{x}\\
 &=\pi^{-s/2}\Gamma(s/2)\sum_{n\geq1}n^{-s}.
\end{align*}
Splitting the integral at $1$, using
$\psi(x)=x^{-1/2}\psi(1/x)+(x^{-1/2}-1)/2$, which follows from
\eqref{eq:theta-transform}, and then substituting $x\mapsto1/x$ in the transformed term gives
\begin{equation}
 \pi^{-s/2}\Gamma(s/2)\zeta(s)
 =\frac{1}{s(s-1)}+
  \int_1^\infty\psi(x)
   \bigl(x^{s/2}+x^{(1-s)/2}\bigr)\frac{\dd x}{x}.
 \label{eq:theta-continuation}
\end{equation}
For $x\geq1$,
$\psi(x)\leq e^{-\pi x}/(1-e^{-3\pi x})$: indeed $n^2\geq1+3(n-1)$ for $n\geq1$.
On every compact set of $s$-values, the integrand and each derivative with respect to $s$ are
therefore dominated by a constant times
$e^{-\pi x}x^C(1+\log^k x)$.  Differentiation under the integral shows that the integral in
\eqref{eq:theta-continuation} is entire.  Multiplication by $s(s-1)/2$ gives
\[
 \xi(s)=\frac12+\frac12s(s-1)
 \int_1^\infty\psi(x)
   \bigl(x^{s/2}+x^{(1-s)/2}\bigr)\frac{\dd x}{x},
\]
which is entire and agrees with \eqref{eq:xi-def} for $\Re s>1$.

It remains to prove the asserted zero-free region.  For $\sigma=\Re s>1$,
$\sum_p|p^{-s}|\leq\sum_{n\geq2}n^{-\sigma}<\infty$, and hence
\[
 \sum_p\sum_{k\geq1}\frac{|p^{-ks}|}{k}<\infty.
\]
Exponentiating this absolutely convergent logarithm gives
$\zeta(s)=\prod_p(1-p^{-s})^{-1}\neq0$.  The Weierstrass product
\[
 \frac1{\Gamma(z)}=ze^{\gamma z}
   \prod_{n\geq1}\left(1+\frac zn\right)e^{-z/n}
\]
is nonzero for $\Re z>0$, so $\Gamma(z)$ is finite and nonzero there.  The remaining factors in
\eqref{eq:xi-def} are nonzero when $\Re s>1$.  Thus $\xi(s)\neq0$ in this half-plane.  Standard
background for these theta, gamma, and zeta identities is given in \cite{Titchmarsh1986}.
\end{proof}

Put
\[
 L(s)=\frac{\xi'(s)}{\xi(s)}.
\]
Lemma~\ref{lem:xi-analytic} makes $L$ holomorphic on $\Re s>1$.

\begin{lemma}[Exact two-step defect]
\label{lem:two-step}
For $\Re s>1$,
\begin{equation}
 L(s+2)-L(s)-\left(\frac1{s+2}+\frac1{s+1}-\frac1{s-1}\right)
 =\mathcal F(s).
 \label{eq:two-step-defect}
\end{equation}
\end{lemma}

\begin{proof}
Logarithmically differentiating \eqref{eq:xi-def} gives
\begin{equation}
 L(s)=\frac1s+\frac1{s-1}-\frac12\log\pi
       +\frac12\Psi(s/2)+\frac{\zeta'(s)}{\zeta(s)},
 \label{eq:xi-logderiv}
\end{equation}
where $\Psi=\Gamma'/\Gamma$.  Absolute convergence of the Euler product gives
\begin{equation}
 \frac{\zeta'(s)}{\zeta(s)}
 =-\sum_{n\geq1}\Lambda(n)n^{-s}.
 \label{eq:zeta-logderiv}
\end{equation}
The gamma recurrence $\Gamma(z+1)=z\Gamma(z)$ implies
$\Psi(z+1)-\Psi(z)=1/z$.  Subtracting \eqref{eq:xi-logderiv} at $s$ from the same formula at
$s+2$, the term $-1/s$ from the rational boundary factor cancels the term $1/s$ from the digamma
recurrence.  The remaining archimedean terms are
\[
 \frac1{s+2}+\frac1{s+1}-\frac1{s-1}.
\]
Finally, \eqref{eq:zeta-logderiv} gives
\[
 \frac{\zeta'(s+2)}{\zeta(s+2)}-
 \frac{\zeta'(s)}{\zeta(s)}
 =\sum_{n\geq1}\Lambda(n)(1-n^{-2})n^{-s}=\mathcal F(s),
\]
which proves \eqref{eq:two-step-defect}.
\end{proof}

\begin{theorem}[Completed-xi hypertranscendence]
\label{thm:xi-ht}
The completed Riemann xi function is hypertranscendental over $\C(s)$.
\end{theorem}

\begin{proof}
Assume that $\xi$ is differentially algebraic.  Restriction to $\Re s>1$ preserves its
differential equation.  Since $\xi$ is nonzero there, Proposition~\ref{prop:da-field} makes
$L=\xi'/\xi$ differentially algebraic.  Lemma~\ref{lem:translation} then makes $L(s+2)$
differentially algebraic, and the rational function in parentheses in
\eqref{eq:two-step-defect} is differentially algebraic.  Proposition~\ref{prop:da-field} and
Lemma~\ref{lem:two-step} would therefore make $\mathcal F$ differentially algebraic.  This
contradicts Proposition~\ref{prop:filtered-ht}.
\end{proof}

\begin{corollary}[Polynomial-gauge invariance of xi]
\label{cor:xi-gauge}
For every $Q\in\C[s]$, both $e^{Q(s)}\xi(s)$ and $e^{-Q(s)}\xi(s)$ are
hypertranscendental over $\C(s)$.
\end{corollary}

\begin{proof}
Apply Lemma~\ref{lem:exp-gauge} and Theorem~\ref{thm:xi-ht}.
\end{proof}

\section{Elimination of the Painlev\'e sigma function}
\label{sec:tau}

Let $V\subset\C$ be open, and let $u,\sigma:V\to\C$ be holomorphic.  Assume the
cross-multiplied tau identity
\begin{equation}
 u\sigma=tu'
 \label{eq:cross-main}
\end{equation}
and the $D_6$ sigma equation
\begin{equation}
 (t\sigma'')^2
 =\bigl(2\sigma'-\theta_0-\theta_\infty\bigr)^2
      (\sigma-t\sigma')
   +4\sigma'(\sigma'-\theta_0)(\sigma'-\theta_\infty)
 \label{eq:sigma-main}
\end{equation}
on $V$.  Write $u_j=u^{(j)}$ and define
\begin{align}
 A&=uu_1+tuu_2-tu_1^2,\label{eq:A-def}\\
 B&=2u(uu_2-u_1^2)
   +t\bigl(u^2u_3-3uu_1u_2+2u_1^3\bigr),\label{eq:B-def}\\
 D&=u_1^2-uu_2.\label{eq:D-def}
\end{align}

\begin{proposition}[Denominator-free tau equation]
\label{prop:tau-equation}
The function $u$ satisfies
\begin{align}
 t^2B^2={}&t^2\bigl(2A-(\theta_0+\theta_\infty)u^2\bigr)^2D
 \notag\\
 &\quad+4A(A-\theta_0u^2)(A-\theta_\infty u^2)
 \label{eq:cleared-tau}
\end{align}
at every point of $V$.  The differential polynomial obtained by moving the right-hand side to
the left is nonzero for every $\theta_0,\theta_\infty\in\C$.
\end{proposition}

\begin{proof}
Differentiate \eqref{eq:cross-main} once:
\[
 u_1\sigma+u\sigma'=u_1+tu_2.
\]
Multiply by $u$ and use $u\sigma=tu_1$.  This gives the first elimination identity
\begin{equation}
 u^2\sigma'=uu_1+tuu_2-tu_1^2=A.
 \label{eq:first-elimination}
\end{equation}
It follows, again without division, that
\begin{align}
 u^2(\sigma-t\sigma')
 &=u(u\sigma)-t(u^2\sigma')
 \notag\\
 &=tuu_1-tA
 =t^2(u_1^2-uu_2)=t^2D.
 \label{eq:difference-elimination}
\end{align}

Differentiate \eqref{eq:first-elimination}:
\[
 2uu_1\sigma'+u^2\sigma''=A'.
\]
After multiplication by $u$ and one use of $u^2\sigma'=A$, this becomes
\begin{equation}
 u^3\sigma''=uA'-2u_1A.
 \label{eq:second-elimination-raw}
\end{equation}
Direct differentiation of \eqref{eq:A-def} gives
\[
 A'=2uu_2+tuu_3-tu_1u_2.
\]
Substitution into \eqref{eq:second-elimination-raw} yields
\begin{align*}
 uA'-2u_1A
 &=2u^2u_2+t u^2u_3-tu u_1u_2
   -2uu_1^2-2tu u_1u_2+2tu_1^3\\
 &=2u(uu_2-u_1^2)
   +t(u^2u_3-3uu_1u_2+2u_1^3)=B.
\end{align*}
Hence
\begin{equation}
 u^3\sigma''=B.
 \label{eq:second-elimination}
\end{equation}

Multiply \eqref{eq:sigma-main} by $u^6$.  Equations
\eqref{eq:first-elimination}--\eqref{eq:second-elimination} give, term by term,
\begin{align*}
 u^6(t\sigma'')^2&=t^2B^2,\\
 u^2(2\sigma'-\theta_0-\theta_\infty)
   &=2A-(\theta_0+\theta_\infty)u^2,\\
 u^2(\sigma-t\sigma')&=t^2D,\\
 u^2(\sigma'-\theta_0)&=A-\theta_0u^2,\\
 u^2(\sigma'-\theta_\infty)&=A-\theta_\infty u^2.
\end{align*}
Their substitution is exactly \eqref{eq:cleared-tau}.

It remains to verify that this equation is genuine.  Regard its left side minus its right side as
a polynomial in the five independent variables $(t,u,u_1,u_2,u_3)$.  At
\[
 (t,u,u_1,u_2,u_3)=(1,1,0,0,1)
\]
one has $A=0$, $B=1$, and $D=0$.  The polynomial therefore takes the value $1$, independently of
$\theta_0$ and $\theta_\infty$.  It is nonzero.
\end{proof}

\begin{remark}
No step in the proof divides by $u$.  Thus Proposition~\ref{prop:tau-equation} is valid at tau
zeros whenever the cross-multiplied tau--sigma chart is holomorphic there.  For an ordinary chart
on which $u\neq0$, \eqref{eq:cross-main} is exactly
$\sigma=t(\log u)'$, so the proposition applies to the usual isomonodromic definition.
\end{remark}

\begin{corollary}
\label{cor:tau-da}
Every regular local $\mathrm{PIII}(D_6)$ tau function is differentially algebraic over $\C(t)$.
\end{corollary}

\begin{proof}
Equation \eqref{eq:cleared-tau} is an algebraic differential equation of order three with
polynomial coefficients in $t$, and Proposition~\ref{prop:tau-equation} proves that its defining
differential polynomial is nonzero.
\end{proof}

\subsection{Why the isomonodromic tau function supplies a regular chart}
\label{subsec:tau-chart}

On a simply connected regular domain of an isomonodromic deformation, the Jimbo--Miwa--Ueno
one-form is closed and the tau function is defined, up to a nonzero constant, by
\[
 \dd\log u=\omega_{\mathrm{JMU}}.
\]
Equivalently, after choosing a base point $t_0$ and normalization $u(t_0)\neq0$,
\[
 u(t)=u(t_0)\exp\left(\int_{t_0}^{t}\omega_{\mathrm{JMU}}\right).
\]
Path independence follows from closedness on the simply connected domain.  The displayed formula
makes $u$ holomorphic and nowhere zero there.  For the Painlev\'e~III Hamiltonian reduction,
$\omega_{\mathrm{JMU}}=H(t)\,\dd t$ and $\sigma=tH$, whence
$u\sigma=tu'$.  Okamoto's Hamiltonian formulation supplies
\eqref{eq:sigma-main}; Miwa's analyticity theorem supplies regular analytic domains
\cite{JMU1981,Miwa1981,Okamoto1981,Okamoto1987}.  Thus every regular isomonodromic point has a
neighborhood satisfying exactly the two identities assumed in
Proposition~\ref{prop:tau-equation}.  No global meromorphic continuation of $u$ is required.

Planat's equation~(5) is \eqref{eq:sigma-main}, with the same parameters
$\theta_0,\theta_\infty$, and his equation~(18) is \eqref{eq:planat-carrier}
\cite{Planat2026}.  Hence the local object used below is not an enlarged surrogate for the
proposed carrier: it is the restriction of that carrier to any regular isomonodromic chart.

\section{The polynomial-gauge no-go theorem}
\label{sec:nogo}

Set
\[
 g(s)=s(1-s),
 \qquad g'(s)=1-2s.
\]
Thus $g'$ has the unique zero $s=\tfrac12$.

\begin{proposition}[Differential algebraicity of a pulled-back carrier]
\label{prop:carrier-da}
Let $U$ be a nonempty disc whose closure avoids $\tfrac12$, and let $(u,\sigma)$ be a regular
local $D_6$ tau chart on $g(U)$.  If an entire function $E$ satisfies
\[
 E(s)=e^{-2s}u(g(s))\qquad(s\in U),
\]
then $E$ is differentially algebraic over $\C(s)$.
\end{proposition}

\begin{proof}
Corollary~\ref{cor:tau-da} gives a nonzero differential equation for $u$ on $g(U)$.  Since
$g'\neq0$ on $U$, Proposition~\ref{prop:pullback} gives a nonzero differential equation over
$\C(s)$ for $u\circ g$ on $U$.  Lemma~\ref{lem:exp-gauge} and
Proposition~\ref{prop:da-field} give a local differential equation for
$e^{-2s}u(g(s))$, hence for $E$.  Lemma~\ref{lem:local-global} extends that equation to all of
$\C$.
\end{proof}

\begin{theorem}[Degree-unrestricted bridge exclusion]
\label{thm:main}
For no $\theta_0,\theta_\infty\in\C$, polynomial $Q\in\C[s]$, function $E:\C\to\C$, nonempty
disc $U$ whose closure avoids $\tfrac12$, and regular local $D_6$ tau chart $(u,\sigma)$ on $g(U)$
do both
identities
\begin{align*}
 E(s)&=e^{-2s}u(g(s)) &&(s\in U),\\
 \xi(s)&=e^{Q(s)}E(s) &&(s\in\C)
\end{align*}
hold.
\end{theorem}

\begin{proof}
Assume that all the displayed data exist.  The global bridge identity gives
\begin{equation}
 E(s)=e^{-Q(s)}\xi(s).
 \label{eq:E-entire}
\end{equation}
Both factors on the right are entire, so $E$ is entire.  Proposition~\ref{prop:carrier-da} now
shows that $E$ is differentially algebraic over $\C(s)$.  Lemma~\ref{lem:exp-gauge} and the bridge
identity then make $\xi=e^QE$ differentially algebraic.  This contradicts
Theorem~\ref{thm:xi-ht}.
\end{proof}

\begin{corollary}[Falsity of the polynomial-gauge bridge]
\label{cor:planat}
There is no choice of fixed $\mathrm{PIII}(D_6)$ monodromy data and no polynomial $Q$ of degree at
most two for which Planat's carrier \eqref{eq:planat-carrier} satisfies
\eqref{eq:planat-bridge}.  More strongly, no polynomial degree works.
\end{corollary}

\begin{proof}
An isomonodromic tau function and its sigma function obey \eqref{eq:tau-sigma-quotient} and
\eqref{eq:sigma-form-intro} on every regular zero-free local chart by
Section~\ref{subsec:tau-chart}.  Choose a regular point $t_0\neq\tfrac14$; this is possible because
a regular chart is open and the excluded set has one point.  The equation $s(1-s)=t_0$ has a
complex solution $s_0$, and $s_0\neq\tfrac12$ because $g(\tfrac12)=\tfrac14$.  Hence
$g'(s_0)=1-2s_0\neq0$.  By continuity, a sufficiently small disc $U$ about $s_0$ has closure
disjoint from $\tfrac12$ and is mapped into the chosen chart.  Equations
\eqref{eq:planat-carrier} and
\eqref{eq:planat-bridge} would therefore supply exactly the data forbidden by
Theorem~\ref{thm:main}.  The theorem places no restriction on $\deg Q$.
\end{proof}

\begin{remark}[Normalization freedom]
Planat also records the tau normalization
$u(t)\mapsto e^{at^2+bt+c}u(t)$.  Under $t=s(1-s)$, its additional exponent is
\begin{equation}
 a[s(1-s)]^2+b[s(1-s)]+c
 =as^4-2as^3+(a-b)s^2+bs+c,
 \label{eq:normalization-polynomial}
\end{equation}
a polynomial of degree at most four.  It can be absorbed into the external polynomial gauge.
The degree-unrestricted form of Theorem~\ref{thm:main} therefore covers this normalization
ambiguity as well as the degree-two gauge stated in Conjecture~3.
\end{remark}

\begin{corollary}[Disjointness of the two classes]
\label{cor:disjointness}
The polynomial-exponential gauge orbit of $\xi$ is disjoint from the class of entire functions
that admit, on one regular quadratic-pullback chart, a fixed-parameter
$\mathrm{PIII}(D_6)$ tau representation.
\end{corollary}

\begin{proof}
Every member of the latter class is differentially algebraic by
Proposition~\ref{prop:carrier-da}.  Every member of the former class is hypertranscendental by
Corollary~\ref{cor:xi-gauge}.  An element cannot belong to both classes by definition.
\end{proof}

\section{Consequences and scope}
\label{sec:scope}

\subsection{Interpretation of the obstruction}

The proof separates arithmetic complexity from isomonodromic dynamics.  On the arithmetic side,
the logarithmic derivative converts the Euler product into prime-power coefficients.  Translation
by two multiplies those coefficients by $1-n^{-2}$, which removes no prime pivot.  Since the
$j$th derivative multiplies the coefficient at $p$ by $(-\log p)^j$, distinct primes provide a
Vandermonde matrix of every size.  Thus no finite differential jet can satisfy an algebraic
relation.  The exact shift identity transfers this obstruction from the filtered prime series to
$\xi$.

On the isomonodromic side, the sigma equation is a finite-order algebraic ordinary differential
equation.  Substitution of $\sigma=t(\log u)'$ might appear to introduce denominators and lose
information at zeros of $u$.  Equations \eqref{eq:first-elimination}--
\eqref{eq:second-elimination} show that multiplication by the exact powers of $u$ clears every
denominator and leaves the genuine relation \eqref{eq:cleared-tau}.  The parameters
$\theta_0,\theta_\infty$ select solutions of this equation but do not alter its
differential-algebraic character.

The contradiction therefore compares two invariant classes, not two approximations.  A polynomial
gauge cannot help because $e^Q$ satisfies $v'-Q'v=0$.  The critical point of $s(1-s)$ cannot help
because a single regular chart away from $s=\tfrac12$ yields a local differential equation, and
entireness propagates that equation by the identity theorem.  Possible tau zeros cannot invalidate
the elimination because no division by $u$ occurs.  Exceptional fixed monodromy data cannot make
the cleared equation vanish identically, as the parameter-independent evaluation in
Proposition~\ref{prop:tau-equation} proves.

\subsection{What the method enables}

The argument gives a screening procedure for exact special-function bridges.  On the proposed
model side, eliminate auxiliary variables until a nonzero algebraic differential equation for the
claimed carrier is obtained; then check that the proposed pullback and gauge preserve differential
algebraicity.  On the target side, seek an exact translation or dilation defect whose Dirichlet
coefficients retain nonzero values on arbitrarily many multiplicatively independent generators.
The prime-specialization and weighted-uniqueness arguments of
Section~\ref{sec:dirichlet} then test whether the target can belong to the same differential class.
This can rule out an entire family of parameter choices at once, before numerical matching or
monodromy reconstruction is attempted.

For the present conjecture, this screening is decisive: changing monodromy data, increasing the
external polynomial degree, moving the local chart, or changing the quoted tau normalization does
not alter the contradiction.  Any viable replacement must change at least one structural input.
Examples include an infinite-dimensional system, a nonalgebraic dependence on the independent
variable, a gauge that is itself hypertranscendental, or a representation that is not an identity
of analytic functions.  These are logical possibilities, not existence claims.

\subsection{Boundary of the result}

The theorem does not exclude every Riemann--Hilbert, canonical-system, or integrable representation
of $\xi$.  It excludes the fixed-parameter $\mathrm{PIII}(D_6)$ tau architecture specified by
\eqref{eq:sigma-form-intro}--\eqref{eq:planat-bridge}, and more generally its polynomial-exponential
gauge orbit on any regular quadratic-pullback chart.  It makes no assertion about approximate
representations, limits of increasing-dimensional systems, or constructions whose coefficients
already contain hypertranscendental data.

No statement about the Riemann hypothesis is used.  The contradiction occurs before zero
localization enters: it is an incompatibility between differential-algebraic types.

\begingroup
\small
\begin{thebibliography}{99}

\bibitem{CashwellEverett1959}
E.~D. Cashwell and C.~J. Everett,
The ring of number-theoretic functions,
\emph{Pacific J. Math.} \textbf{9} (1959), no.~4, 975--985,
\href{https://doi.org/10.2140/pjm.1959.9.975}{doi:10.2140/pjm.1959.9.975}.

\bibitem{HanLiu2021}
Q.~Han and J.~Liu,
Algebraic differential independence regarding the Riemann $\zeta$-function and the Euler
$\Gamma$-function,
\emph{J. Number Theory} \textbf{221} (2021), 109--121,
\href{https://doi.org/10.1016/j.jnt.2019.12.006}{doi:10.1016/j.jnt.2019.12.006}.

\bibitem{HardyRiesz1915}
G.~H. Hardy and M.~Riesz,
\emph{The General Theory of Dirichlet's Series},
Cambridge Tracts in Mathematics and Mathematical Physics, vol.~18,
Cambridge University Press, Cambridge, 1915.

\bibitem{JMU1981}
M.~Jimbo, T.~Miwa, and K.~Ueno,
Monodromy preserving deformation of linear ordinary differential equations with rational
coefficients: I. General theory and $\tau$-function,
\emph{Physica D} \textbf{2} (1981), no.~2, 306--352,
\href{https://doi.org/10.1016/0167-2789(81)90013-0}{doi:10.1016/0167-2789(81)90013-0}.

\bibitem{Kolchin1973}
E.~R. Kolchin,
\emph{Differential Algebra and Algebraic Groups},
Pure and Applied Mathematics, vol.~54, Academic Press, New York, 1973.

\bibitem{Markus2007}
L.~Markus,
Differential independence of $\Gamma$ and $\zeta$,
\emph{J. Dynam. Differential Equations} \textbf{19} (2007), no.~1, 133--154,
\href{https://doi.org/10.1007/s10884-006-9034-1}{doi:10.1007/s10884-006-9034-1}.

\bibitem{Miwa1981}
T.~Miwa,
Painlev\'e property of monodromy preserving deformation equations and the analyticity of
$\tau$ functions,
\emph{Publ. Res. Inst. Math. Sci.} \textbf{17} (1981), no.~2, 703--721,
\href{https://doi.org/10.2977/prims/1195185270}{doi:10.2977/prims/1195185270}.

\bibitem{Nagoshi2016}
H.~Nagoshi,
Hypertranscendence of $L$-functions for $\mathrm{GL}_m(\mathbb A_{\mathbb Q})$,
\emph{Bull. Aust. Math. Soc.} \textbf{93} (2016), no.~3, 388--399,
\href{https://doi.org/10.1017/S000497271500129X}{doi:10.1017/S000497271500129X}.

\bibitem{Okamoto1981}
K.~Okamoto,
On the $\tau$-function of the Painlev\'e equations,
\emph{Physica D} \textbf{2} (1981), no.~3, 525--535,
\href{https://doi.org/10.1016/0167-2789(81)90026-9}{doi:10.1016/0167-2789(81)90026-9}.

\bibitem{Okamoto1987}
K.~Okamoto,
Studies on the Painlev\'e equations. IV. Third Painlev\'e equation $P_{\mathrm{III}}$,
\emph{Funkcialaj Ekvacioj} \textbf{30} (1987), 305--332.

\bibitem{Ostrowski1920}
A.~Ostrowski,
\"Uber Dirichletsche Reihen und algebraische Differentialgleichungen,
\emph{Math. Z.} \textbf{8} (1920), 241--298.

\bibitem{Planat2026}
M.~Planat,
Wild character varieties, Painlev\'e $\mathrm{III}_{D_6}$, the isomonodromic cosine, and a complete
positivity framework toward the Riemann hypothesis,
Preprints 2026, 202604.0442,
\href{https://doi.org/10.20944/preprints202604.0442.v1}
  {\nolinkurl{doi:10.20944/preprints202604.0442.v1}}.

\bibitem{Reich1984}
A.~Reich,
On hypertranscendental functions,
in \emph{Topics in Classical Number Theory}, Vol.~II,
Colloquia Mathematica Societatis J\'anos Bolyai, vol.~34,
North-Holland, Amsterdam, 1984, pp.~1349--1369.

\bibitem{Ritt1950}
J.~F. Ritt,
\emph{Differential Algebra},
American Mathematical Society Colloquium Publications, vol.~33,
American Mathematical Society, New York, 1950.

\bibitem{Titchmarsh1986}
E.~C. Titchmarsh,
\emph{The Theory of the Riemann Zeta-Function}, 2nd ed., revised by D.~R. Heath-Brown,
Clarendon Press, Oxford, 1986.

\bibitem{WangYao2024}
Q.~Wang and X.~Yao,
Differential and difference independence of $\zeta$ and $\Gamma$,
\emph{J. Differential Equations} \textbf{412} (2024), 797--807,
\href{https://doi.org/10.1016/j.jde.2024.08.068}{doi:10.1016/j.jde.2024.08.068}.

\end{thebibliography}
\endgroup

\end{document}
